\section{引言：AI 与生物学的交汇}

本讲是 MIT 6.S191 (Introduction to Deep Learning) 2025 春季学期的最后一讲，由 Microsoft Research (MSR) 研究员 Ava Amini 作为特邀嘉宾授课。Amini 拥有横跨湿实验室生物学与计算科学的双重背景——本科和博士期间深耕细胞实验与动物模型研究，后期将对生物学的热情与计算方法结合，致力于以 AI 手段推动生命科学的前沿发展。

\begin{knowledgebox}{Microsoft Research (MSR) 的使命}
MSR 的核心使命是"推进科学与技术前沿，造福人类" (advance the frontiers of science and technology to benefit humanity)。其研究以三个层次的视角展开：(1) 基础性、学术驱动的前沿研究；(2) 负责任且符合伦理的技术部署；(3) 对人类健康与自然科学发现产生积极影响。MSR 内部设有专门的 Biomedical Machine Learning (BioML) 团队，Amini 正是该团队的核心成员。
\end{knowledgebox}

Amini 指出，生物学本质上是一个在纳米尺度运作的极端复杂系统。即便是单个细胞膜边界在 100 万倍放大下的科学可视化，也展现出惊人的结构丰富度——大量的蛋白质分子在其中承担各种功能。这种复杂性既是 AI 面临的巨大挑战，也是前所未有的机遇。

在 6.S191 课程中，学生已经学习了 \textbf{预测建模} (predictive modeling) 和 \textbf{生成建模} (generative modeling) 两大范式。当这一框架被应用于生物领域时，核心问题依然是预测与生成，但数据模态发生了根本性转变：

\begin{itemize}
    \item 不再是自然语言文本，而是单个 \textbf{生物分子的序列语言}
    \item 不再是人脸或汽车图像，而是 \textbf{细胞表征}、\textbf{组织图像}、\textbf{患者数据}
    \item 预测目标变为：生物分子的功能是什么？细胞对药物如何响应？患者的临床疗效如何？
    \item 生成目标变为：给定目标功能，能否设计出实现该功能的生物分子？
\end{itemize}

\begin{importantbox}{AI for Biology 的核心双向范式}
\textbf{正向预测}：给定生物分子 $\rightarrow$ 预测其功能（分类、回归）。\\
\textbf{反向设计}：给定目标功能 $\rightarrow$ 生成满足该功能的生物分子（生成模型）。\\
这两个方向互为补充，共同构成 AI 赋能生物学的完整闭环。关键在于，AI 的预测和设计最终都需要通过\textbf{实验验证}与自然世界对接——这是 AI for Biology 区别于纯计算领域的核心特征。
\end{importantbox}

本讲的核心内容聚焦于\textbf{生成式 AI 在蛋白质设计中的应用}，特别是如何将 Diffusion Model（扩散模型）这一强大的生成建模框架从图像领域迁移到蛋白质序列的离散数据域，并介绍了 Amini 团队开发的 \textbf{EvoDiff} 模型。

\subsection{本章小结}

AI for Biology 是 AI 最具潜力的应用方向之一。该领域的核心范式包括预测建模和生成建模两大支柱，而其独特之处在于必须通过实验与自然世界对接。本讲聚焦于利用生成式 AI（特别是 Diffusion Model）来设计具有目标功能的蛋白质分子。


